\section*{Abstract}
\addcontentsline{toc}{section}{Abstract}

This thesis asks two questions: whether Heterogeneous Autoregressive (HAR) quantile regression produces better out-of-sample distributional forecasts of cocoa futures volatility than established benchmarks, and whether augmenting the model with the El~Ni\~no--Southern Oscillation (ENSO) climate index improves long-horizon upper-tail forecasts. Daily Open-High-Low-Close prices for ICE London cocoa front-month futures spanning January~1985 to December~2025 (10,376 trading days) are used to construct a Yang-Zhang volatility proxy, validated against realized volatility from approximately one year of five-minute intraday data. Forecasts at six horizons (1~day to 12~months) are evaluated by quantile loss and QLIKE against five competing models, with Model Confidence Sets used to assess statistical significance.

On the first question, the value of direct quantile estimation is concentrated at short horizons. HAR-QR delivers the lowest upper-tail ($\tau = 0.95$) quantile loss from one day through three months, and at one day and one week it is the sole top-ranked model in a 90\% Model Confidence Set that excludes historical volatility, GARCH(1,1), and the HAR-OLS pseudo-quantile. The ranking then reverses: the HAR-OLS pseudo-quantile overtakes HAR-QR at six months, and GARCH(1,1) is the best upper-tail forecaster at twelve months, 29\% below HAR-QR in quantile loss, though these long-horizon differences are not statistically significant. For mean forecasts, HAR-OLS dominates at every horizon under QLIKE, MSE, and MAE. Upper-quantile coverage is near the nominal $5\%$ at one day and one week but deteriorates as the horizon lengthens, reaching 14.6\% at twelve months.

On the second question, the Ni\~no~3.4 index worsens upper-tail quantile loss at every horizon from one to six months and improves it by only $3.3\%$ at twelve months. Three robustness diagnostics dissolve even this residual gain: it is confined entirely to the 2022--2024 cocoa supply-crisis window---in the pre-crisis subsample ENSO augmentation \emph{raises} twelve-month upper-tail loss by 16\%---it follows a lag structure that contradicts the agronomic transmission hypothesis, and it does not replicate on Arabica coffee or raw sugar. ENSO therefore yields no robust forecast improvement at any horizon; its apparent long-horizon value is a regime-conditional, sample-specific artefact rather than a replicable structural climate channel.

The thesis makes three contributions: the first application of HAR-QR to an agricultural commodity, the first evaluation at horizons extending to twelve months, and the first empirical test of an ENSO-augmented HAR-QR specification for cocoa.

\vspace{0.5cm}
\noindent\textbf{Keywords:} cocoa futures, volatility forecasting, HAR model, quantile regression, ENSO, long-horizon forecasting, Yang-Zhang estimator
